\documentclass[10pt,a4paper]{amsart}
\usepackage[margin=19mm]{geometry}
\usepackage{amsmath,amssymb,mathtools}
\usepackage[scr=rsfs,cal=euler]{mathalfa}
\usepackage{xfrac}
\usepackage[unicode,hidelinks,bookmarks=false]{hyperref}
\newcommand{\ie}{i.e.\ }
\newcommand{\cf}{cf.\ }
\newcommand{\resp}{respectively\ }
\newcommand{\Subsection}{\subsection}
\providecommand{\nobreakdash}{\nobreak\hbox{-}}
\setlength{\emergencystretch}{2em}
\multlinegap0pt
\usepackage{amssymb}
\usepackage[shortlabels]{enumitem}
\usepackage{xcolor}
\newtheorem{definition}{Definition}
\newtheorem{lemma}{Lemma}
\newtheorem{remark}{Remark}
\newtheorem{overpic}{Overpic}
\title{Complementary legs and symplectic rational balls}
\author{John B. Etnyre, Burak Ozbagci,, B\"ulent Tosun}
\date{Source: arXiv:2505.04513v2 (CC BY 4.0)}
\begin{document}
\maketitle

\noindent\emph{Source and licence.} Complementary legs and symplectic rational balls. Version arXiv:2505.04513v2 (CC BY 4.0), distributed by arXiv under the Creative Commons Attribution 4.0 International licence, http://creativecommons.org/licenses/by/4.0/, as recorded in the arXiv licence metadata for this exact version and shown on the abstract page https://arxiv.org/abs/2505.04513v2. Full licence notice: \texttt{licenses/arxiv-2505.04513-CC-BY-4.0.txt}.\par\smallskip

\noindent\textit{Editorial scope.} Counted unit: the article's lemma inverse (source0.tex lines 886--898). The article's complete original proof of it is delivered with it (source0.tex lines 1324--1378, one \texttt{proof} environment of 55 lines, which the article places away from the statement). No mathematics is written, repaired or re-derived: the statement and the proof are the article's own lines, in the article's own notation, and no retained line is substituted or re-flowed.  Retained uncounted as the source-local context the proof needs: lines 788--794; lines 798--800; lines 845--846; lines 869--883; lines 900--917; lines 919--920.  Nothing was omitted from any retained span, so the package carries no omission record.  No retained line contains a citation, so the wrapper carries no bibliography and no bibliography entry.  No retained line contains a figure environment, an image-inclusion command or drawing code, so no image is delivered and the package declares no asset dependency.  Editorial formatting only, in addition: an AMS \texttt{amsart} preamble that restates the article's own declarations and macros used by the retained lines, namely the environments \texttt{definition}, \texttt{lemma}, \texttt{remark}, \texttt{overpic} on separate counters, together with the standard packages those lines need.  The retained environments print in this document as Definition 1, Lemma 1, Remark 1, Overpic 1 in order of appearance, whereas the article prints the same items with the numbering of its own full document.  The complete native source of the article as distributed in the arXiv e-print for this exact version is bundled unchanged as \texttt{sources/177754/source0.tex}.
\begin{definition} \label{def: det} For any integer $n>0$ and any sequence of integers $m_1, m_2, \ldots, m_n$ with $m_i \geq 2$, we set $$M=M(m_1, m_2, \ldots, m_n)= \left[\begin{array}{cccc}
 -m_1 & 1 & &   \\
1 & -m_2 & \ddots & \\
 & \ddots & \ddots & 1 \\
 & & 1 & -m_n 
\end{array}\right].$$  For any $1 \leq i \leq n$, we define $u^M_i$ as the absolute value of the determinant of the top-left $i \times i$ submatrix of $M$, and we set  $u^M_0=1$.  Let ${\bf u}^M=[u^M_0 \; u^M_1 \; \cdots \; u^M_{n-1}]^T \in \mathbb{R}^n$. Similarly, for any $1 \leq j \leq n$, we define $v^M_j$ as the absolute value of the determinant of the bottom-right $j \times j$ submatrix of $M$, and we set $v^M_0=1$. Let ${\bf v}^M=[v^M_{n-1} \; v^M_{n-2} \; \cdots \; v^M_0]^T \in \mathbb{R}^n$. 
\end{definition}

\begin{lemma} \label{lem: firstrow} 
Using the notation as in Definition~\ref{def: det}, let ${s}/{t}=[m_1, m_2, \ldots, m_n]$ for some relatively prime integers $0 < t < s$. Then $\det M = (-1)^n s$,   the first column of $M^{-1}$ is given by the vector  $-\dfrac{1}{s}   \; {\bf v}^M$ and its dot product with the vector $[m_1-2 \; m_2-2 \; \cdots \; m_n-2]^T \in  \mathbb{R}^n$ is equal to $-1+ \dfrac{1+t}{s}$. Moreover,  the last column of $M^{-1}$ is given by the vector  $-\dfrac{1}{s}   \; {\bf u}^M$. Furthermore, the first entry of  ${\bf v}^M$ is $t$. 
\end{lemma} 

\begin{remark} \label{rem: specialcase} In Definition~\ref{def: det}, we assumed that $m_i \geq 2$ for all $1 \leq i \leq n$. Suppose that we relax this condition so that $m_1=1$ and $m_i \geq 2$ for all $2 \leq i \leq n$, and let $M=M(1, m_2, \ldots, m_n)$, and the vectors ${\bf u}^M$ and ${\bf v}^M$ be defined as in Definition~\ref{def: det}. Then  Lemma~\ref{lem: firstrow} still (partially) holds as follows. We first observe that  $[1, m_2, \ldots, m_n]=\dfrac{s}{t}$ for some uniquely determined co-prime integers $0 < s < t $ (as opposed to $0 < t < s$). Then   $\det M = (-1)^n s$, the first and last columns of $M^{-1}$ are given by   $-(1/s)   \; {\bf v}^M$ and   $-(1/s)   \; {\bf u}^M$, respectively.  
\end{remark}

\begin{figure}[htb]{\small
\begin{overpic}%[grid,tics=10] 
{matrix}
{\Large
\put(30, 155){$A$}
\put(95, 90){$B$}
\put(158, 28){$C$}
\put(77, 128){$1$}
\put(63, 110){$1$}
\put(136, 110){$1$}
\put(76, 52){$1$}}
\end{overpic}}
\caption{Intersection matrix $Q_{{\bf a^1}, {\bf a^2}}$ for $X_{{\bf a^1}, {\bf a^2}}$.}
\label{fig: matrix}
\end{figure}

\begin{figure}[htb]{\small
\begin{overpic}%[grid,tics=10] 
{matrix}
{\Large
\put(10, 154){$A^{-1}+G$}
\put(91, 91){$\widetilde{B}^{-1}$}
\put(138, 28){$C^{-1}+H$}
\put(95, 154){$D$}
\put(27, 91){$D^T$}
\put(157, 91){$F$}
\put(95, 28){$F^T$}
\put(157, 154){$E$}
\put(27, 28){$E^T$}
}
\end{overpic}}
\caption{The matrix $Q^{-1}_{{\bf a^1}, {\bf a^2}}$.}
\label{fig: inversematrix}
\end{figure}

\begin{remark} \label{rem: row} Note that in the definition of the matrix $E$, for example, ${\bf u}^A ( {\bf v}^C)^T$ is the $n_1 \times n_3$ matrix obtained by  multiplying the $n_1 \times 1$ column matrix ${\bf u}^A$ with the $1 \times n_3$ row matrix $({\bf v}^C)^T$. Since the first component of ${\bf u}^A$ is equal to $1$, by definition, this means that the first row of the matrix $E$ is given by $ -\dfrac{q}{(\widetilde{p})^2(p-q)}  ({\bf v}^C)^T$ and the $j$th row of $E$ is obtained by multiplying its first row by the $j$th component of ${\bf u}^A$.  A similar discussion applies to the matrices $D,  F,  G, H$ and $D^T, E^T, F^T$ as well.  
\end{remark} 

\begin{lemma} \label{lem: inverse} The inverse $Q^{-1}_{{\bf a^1}, {\bf a^2}}$ can be described as in Figure~\ref{fig: inversematrix}, where the block matrices 
$\widetilde{B}$, $D$, $E$, $F$, $G$, $H$  in the figure  are defined as follows: 
\[
G= -\dfrac{q}{(\widetilde{p})^2(p-q)}  {\bf u}^A ( {\bf u}^A)^T,  \;\;\; 
D = -\dfrac{1}{\widetilde{p} (p-q)}   {\bf u}^A ({\bf v}^B)^T,  \;\;\; 
E = -\dfrac{q}{(\widetilde{p})^2(p-q)}   {\bf u}^A  ({\bf v}^C)^T
\] 
\[ 
\widetilde{B}=M(a^2_0-1, a^2_1, \ldots, a^2_{n_2}), \;\;\; 
F=  -\dfrac{1}{\widetilde{p} (p-q)}   {\bf v}^B ({\bf v}^C)^T, \;\;\; 
H= -\dfrac{q}{(\widetilde{p})^2(p-q)}   {\bf v}^C ({\bf v}^C)^T .
\] 
\end{lemma} 

 \begin{proof} [Proof of Lemma~\ref{lem: inverse}]  
We have $\widetilde{p}/\widetilde{q}=[a^1_1, a^1_2, \ldots, a^1_{n_1}]$,  $p/q=[a^2_0, a^2_1, \ldots, a^2_{n_2}]$, and $\widetilde{p}/(\widetilde{p}-\widetilde{q})=[a^3_1, a^3_2, \ldots, a^3_{n_3}]$, with all continued fraction coefficients $a^j_i$ being greater than or equal to $2$, by definition. Note that we also have, $\widetilde{p}/(\widetilde{q})^*=[a^1_{n_1}, a^1_{n_2}, \ldots, a^1_1]$, where $(\widetilde{q})^*$ is the inverse of $\widetilde{q}$ mod $\widetilde{p}$,    and $(p-q)/q=[a^2_0-1, a^2_1, \ldots, a^2_{n_2}]$ (see Remark~\ref{rem: specialcase}, when $a^2_0=2$). 

We observe the fact that ${\bf v}^B={\bf v}^{\widetilde{B}}$, which immediately follows from the definitions of the matrices $B$ and $\widetilde{B}$.  By Lemma~\ref{lem: firstrow},  we know that the last column of $A^{-1}$ is  $-(1/\widetilde{p})\;  {\bf u}^A$,   the first column of  $(\widetilde{B})^{-1}$ is  $-(1/(p-q)) \; {\bf v}^B$ (and hence the first row of 
$(\widetilde{B})^{-1}$ is given by $-(1/(p-q)) \; ({\bf v}^B)^T$) and  the first column of  $C^{-1}$ is  $-(1/\widetilde{p})\;  {\bf v}^C$.

The proof of the lemma will be achieved by showing that the product of the matrix $Q_{{\bf a^1}, {\bf a^2}}$, depicted in Figure~\ref{fig: matrix}, and the matrix depicted in Figure~\ref{fig: inversematrix} is equal to the identity matrix. Towards that goal, we first define an auxiliary matrix $\widetilde{Q}_{{\bf a^1}, {\bf a^2}}$ by deleting $A^{-1}$ and $C^{-1}$ from the matrix depicted in Figure~\ref{fig: inversematrix}, and calculate the product of $Q_{{\bf a^1}, {\bf a^2}}$ and  $\widetilde{Q}_{{\bf a^1}, {\bf a^2}}$. 

Let $\widetilde{Q}_A$ be the submatrix of $\widetilde{Q}_{{\bf a^1}, {\bf a^2}}$ obtained by juxtaposition of the blocks $G, D$, and $E$, let $\widetilde{Q}_{\widetilde{B}}$ be the submatrix of  $\widetilde{Q}_{{\bf a^1}, {\bf a^2}}$  obtained by juxtaposition of the blocks  $D^T, (\widetilde{B})^{-1}$ and $F$ and let $\widetilde{Q}_C$ be the submatrix of  $\widetilde{Q}_{{\bf a^1}, {\bf a^2}}$ obtained by juxtaposition of the blocks  $E^T, F^T$, and $H$.

By Remark~\ref{rem: row}, we know that each column of $\widetilde{Q}_A$ is a multiple of ${\bf u}^A$ (because  the definition of $G, D, E$ involves ${\bf u}^A$), that each column of $\widetilde{Q}_C$ is a multiple of ${\bf v}^C$ (because  the definition of $E^T, F^T, H$  involves ${\bf v}^C$) and that the columns of $\widetilde{Q}_{\widetilde{B}}$ belonging to the blocks $D^T$ and  $F$  is a multiple of ${\bf v}^B$ (because the definition of $D^T$ and  $F$  involves ${\bf v}^B$). 
Note that the first column of the block $(\widetilde{B})^{-1}$ is also a multiple of ${\bf v}^B$.   

Now we calculate the dot product of each row of $Q_{{\bf a^1}, {\bf a^2}}$ with each column of  $\widetilde{Q}_{{\bf a^1}, {\bf a^2}}$ case by case below. Note that for each dot product there are at most four terms to calculate. So, the calculations are not very difficult but tedious. 

For $1 \leq i \leq n_1-1$, and $1 \leq j \leq n_1+n_2+n_3+1$, the dot product of the $i$th row of $Q_{{\bf a^1}, {\bf a^2}}$  and the $j$th  column of $\widetilde{Q}_{{\bf a^1}, {\bf a^2}}$ is the same as the dot product of the $i$th row of $A$ and the  $j$th  column of $\widetilde{Q}_A$. Hence this is the dot product of the $i$th row of $A$ and a multiple of the last column of $A^{-1}$ since the last column of $A^{-1}$ is a multiple of ${\bf u}^A$ and each column of $\widetilde{Q}_A$ is a multiple of ${\bf u}^A$. Therefore this dot product is zero, because we are indeed taking the dot product of any of the first $n_1-1$ rows of $A$ with a multiple of the last column of $A^{-1}$.  

We want to show  that for $1 \leq j \leq n_1+n_2+n_3+1$, the dot product of the $n_1$th row of $Q_{{\bf a^1}, {\bf a^2}}$ (the last row of $\widetilde{Q}_A$) and the $j$th  column of $\widetilde{Q}_{{\bf a^1}, {\bf a^2}}$ is zero. To see this,  we first claim that the $(n_1+1)$st row of  $\widetilde{Q}_{{\bf a^1}, {\bf a^2}}$ (the first row of $\widetilde{Q}_{\widetilde{B}}$) is $(\det A=\widetilde{p})$ multiple of the first row of $\widetilde{Q}_{{\bf a^1}, {\bf a^2}}$. To prove the claim we observe that the first row of $(\widetilde{B})^{-1}$ is $\frac{-1}{p-q} ({\bf v}^B)^T=\frac{-1}{p-q} ({\bf v}^{\widetilde{B}})^T$ and the first entry of ${\bf v}^B$ is $q$. Then the  claim follows by simply comparing the first rows of  $D^T, (\widetilde{B})^{-1}, F$   with the first rows of $G, D, E,$ respectively as follows. The first row of $$D^T = -\dfrac{1}{\widetilde{p} (p-q)}   {\bf v}^B ({\bf u}^A)^T \; \mbox{is given by} \;    -\dfrac{q}{\widetilde{p} (p-q)} ({\bf u}^A)^T$$ and the first row of $$G= -\dfrac{q}{(\widetilde{p})^2(p-q)}  {\bf u}^A ( {\bf u}^A)^T \; \mbox{is given by} \;   -\dfrac{q}{(\widetilde{p})^2(p-q)} ( {\bf u}^A)^T.$$ As pointed out above, the first row of 
$$(\widetilde{B})^{-1}  \; \mbox{is given by} \; \frac{-1}{p-q} ({\bf v}^B)^T$$ and the first row of $$D = -\dfrac{1}{\widetilde{p} (p-q)}  {\bf u}^A ({\bf v}^B)^T   \; \mbox{is given by} \;    -\dfrac{1}{\widetilde{p} (p-q)} ({\bf v}^B)^T.$$ Similarly,  the first row of $$E = -\dfrac{q}{(\widetilde{p})^2(p-q)}   {\bf u}^A  ({\bf v}^C)^T \; \mbox{is given by} \;    -\dfrac{q}{(\widetilde{p})^2(p-q)} ({\bf v}^C)^T$$ and the first row of $$F=  -\dfrac{1}{\widetilde{p} (p-q)}   {\bf v}^B ({\bf v}^C)^T \; \mbox{is given by} \;   -\dfrac{q}{\widetilde{p} (p-q)} ({\bf v}^C)^T,$$ which finishes the proof of our claim. Note that the $n_1$th row of $Q_{{\bf a^1}, {\bf a^2}}$  can be seen as the  juxtaposition of the last row of $A$ and the row vector $[1, 0, \dots, 0] \in \mathbb{R}^{n_2+n_3+1}$. So, in the dot product at hand we only have to consider the first $(n_1+1)$st entries  in the $j$th column of $\widetilde{Q}_{{\bf a^1}, {\bf a^2}}$. Our analysis above therefore implies that for $1 \leq j \leq n_1+n_2+n_3+1$, the dot product of the $n_1$th row of $Q_{{\bf a^1}, {\bf a^2}}$ and the $j$th column of $\widetilde{Q}_{{\bf a^1}, {\bf a^2}}$ is equal to a multiple of $$(\mbox{last row of}\; A) \cdot {\bf u}^A + \det A$$ but ${\bf u}^A$ is $-\det A$ times the last column of $A^{-1}$, and therefore the dot product is a multiple of 
$$-\det A (\mbox{last row of}\; A) \cdot (\mbox{last column of}\; A^{-1} )+ \det A=0$$
 in this case as well.  

Now we turn our attention to the last $n_3-1$  rows of $Q_{{\bf a^1}, {\bf a^2}}$. For $$n_1+n_2+3 \leq i \leq n_1+n_2+n_3+1\; \mbox{and}\; 1 \leq j \leq n_1+n_2+n_3+1,$$ the dot product of the $i$th row of $Q_{{\bf a^1}, {\bf a^2}}$ and $j$th  column of $\widetilde{Q}_{{\bf a^1}, {\bf a^2}}$ is the same as the dot product of the $(i-n_1-n_2-1)$st  row of $C$ and the  $j$th  column of $\widetilde{Q}_C$. Hence this is the dot product of the $(i-n_1-n_2-1)$st row of $C$ and a multiple of the first column of $C^{-1}$ since the first column of $C^{-1}$ is a multiple of ${\bf v}^C$ and  each column of  $\widetilde{Q}_C$ is a multiple of ${\bf v}^C$.  Therefore this dot product is zero,  because we are indeed taking the dot product of any of the last  $n_3-1$ rows of $C$ with a multiple of the first column of $C^{-1}$. 

For  $1 \leq j \leq n_1+n_2+n_3+1$, the dot product of the $(n_1+n_2+2)$nd row of $Q_{{\bf a^1}, {\bf a^2}}$ and the $j$th  column of $\widetilde{Q}_{{\bf a^1}, {\bf a^2}}$ is zero because of the fact that the $(n_1+1)$st row of  $\widetilde{Q}_{{\bf a^1}, {\bf a^2}}$ (the first row of $\widetilde{Q}_{\widetilde{B}}$) is $(\det A =\widetilde{p}=\det C)$ multiple of the first row of $\widetilde{Q}_{{\bf a^1}, {\bf a^2}}$, as we showed above.   Hence, the dot product at hand is equal to a multiple of $$(\mbox{first row of}\; C) \cdot {\bf v}^C + \det C$$ but ${\bf v}^C$ is $-\det C$ times the first column of $C^{-1}$, and therefore the dot product is zero in this case as well. 

Now we turn our attention to the middle rows of  $Q_{{\bf a^1}, {\bf a^2}}$. For $n_1+2 \leq i \leq n_1+n_2+1$ and ($1 \leq j \leq n_1+1$ and $n_1+n_2+2 \leq j \leq n_1+n_2+n_3+1$) the dot product of the $i$th row of $Q_{{\bf a^1}, {\bf a^2}}$ and the $j$th  column of $\widetilde{Q}_{{\bf a^1}, {\bf a^2}}$ is the same as the dot product of the $(i-n_1)$th row of $B$ and the  $j$th  column of  $\widetilde{Q}_{\widetilde{B}}$. Hence this is the dot product of the $(i-n_1)$th row of $B$ and a multiple of the first column of $(\widetilde{B})^{-1}$ since the first column of $(\widetilde{B})^{-1}$ is a multiple of ${\bf v}^{\widetilde{B}}={\bf v}^B$ and every column except the ones enumerated from $n_1+2$ to $n_1+n_2+1$ of  $\widetilde{Q}_{\widetilde{B}}$ is a multiple of ${\bf v}^B$. Therefore this dot product is zero, because we are taking the dot product of any of the last $n_2$ rows of $B$ (same as the last $n_2$ rows of  $\widetilde{B}$) with a multiple of the first column of $(\widetilde{B})^{-1}$. 

For $n_1+2 \leq i \leq n_1+n_2+1$ and $n_1+2 \leq j \leq n_1+n_2+1$,  the dot product of the $i$th row of $Q_{{\bf a^1}, {\bf a^2}}$  and the $j$th  column of $\widetilde{Q}_{{\bf a^1}, {\bf a^2}}$ gives a $n_2 \times n_2$ identity matrix because these dot products are the same as the dot products of last $n_2$  rows of $\widetilde{B}$ and the last $n_2$ columns of  $(\widetilde{B})^{-1}$. 

In order to move forward with our analysis, we prove a general fact that will be used below. Since each column of $\widetilde{Q}_C$ is a multiple of ${\bf v}^C$, and the first entry of ${\bf v}^C$ is $(\widetilde{p}-\widetilde{q})$, while its last entry is $1$, we see that the first row of $\widetilde{Q}_C$ is $(\widetilde{p}-\widetilde{q})$ multiple of its last row. Recall that we showed above that the first  row of  $\widetilde{Q}_{\widetilde{B}}$ is $\widetilde{p}$ multiple of the first row of $\widetilde{Q}_A$. Since each column of $\widetilde{Q}_A$ is a multiple of ${\bf u}^A$, and the first entry of ${\bf u}^A$ is $1$, while its last entry $\widetilde{q}$, we see that the last row of $\widetilde{Q}_A$ is $\widetilde{q}$ multiple of its first row. By simply comparing the definitions of the block matrices involved, we also observe that the first row of $\widetilde{Q}_A$ is the same as the last row $\widetilde{Q}_C$. Putting all these observations together, we conclude that 
\begin{equation}\label{eq: abcrows}
 \mbox{last row of}\;  \widetilde{Q}_A +  \mbox{first row of}\;  \widetilde{Q}_C = \mbox{first row of}\;  \widetilde{Q}_{\widetilde{B}} 
\end{equation}
which is a key result for the rest of our proof. 

Next we want to see that the dot product of the $(n_1+1)$st row of $Q_{{\bf a^1}, {\bf a^2}}$ and the  $(n_1+1)$st  column of $\widetilde{Q}_{{\bf a^1}, {\bf a^2}}$ is $1$. The equality~(\ref{eq: abcrows}) can be rephrased as 
\begin{equation}\label{eq: abcrows2}
(n_1)\mbox{th row} + (n_1+n_2+2)\mbox{nd row} = (n_1+1)\mbox{th row} 
\end{equation}
for the matrix $Q_{{\bf a^1}, {\bf a^2}}$. Note that  in the  $(n_1+1)$st  row of $Q_{{\bf a^1}, {\bf a^2}}$, the nonzero terms appear on the $n_1$th,  $(n_1+1)$st , $(n_1+2)$nd and $(n_1+n_2+2)$nd entries as $1,-a_0^2,1,1$. Because of the equality~(\ref{eq: abcrows2}),   the dot product at hand is exactly the dot product of the first row of $\widetilde{B}$ (which is the vector $[-(a_0^2-1) \; 1 \; 0 \cdots 0] \in \mathbb{R}^{n_2+1}$ and the first column of $(\widetilde{B})^{-1}$, which is equal to $1$. 

What is left to consider is the dot product of the  $(n_1+1)$st  row of $Q_{{\bf a^1}, {\bf a^2}}$ and the  $j$th column of  $\widetilde{Q}_{{\bf a^1}, {\bf a^2}}$ for  $j \neq n_1+1$. If we take $1 \leq j \leq n_1$, the dot product of the $(n_1+1)$st row of $Q_{{\bf a^1}, {\bf a^2}}$  and the $j$th column of $\widetilde{Q}_{{\bf a^1}, {\bf a^2}}$ is not zero. The computation here is very similar to the case $j=n_1+1$ we discussed in the paragraph above. In fact the dot product here is equal to some multiples (for $1 \leq j \leq n_1$) of $1$ (that we  computed in the previous paragraph), and as a matter of fact  these dot products are exactly given by the row vector $(1/\widetilde{p})\; ({\bf u}^A)^T$. Therefore,  when we compute the dot product  of the $(n_1+1)$st  row of $Q_{{\bf a^1}, {\bf a^2}}$  and 
the $j$th column of claimed $Q_{{\bf a^1}, {\bf a^2}}^{-1}$ (where we have to take into account $A^{-1}$ at the top left block) we just have to add $-(1/\widetilde{p})\; ({\bf u}^A)^T$ (the last row of $A^{-1}$) and hence the dot product will be zero.  

Similarly, if we take $n_1+n_2+2 \leq j \leq n_1+n_2+n_3+1$, the dot product of the $(n_1+1)$st row of $Q_{{\bf a^1}, {\bf a^2}}$  and the $j$th column of $\widetilde{Q}_{{\bf a^1}, {\bf a^2}}$ is not zero and in fact equal to some multiples of $1$ we computed above. These multiples are exactly given by the row vector $(1/\widetilde{p})\; ({\bf v}^C)^T$. Hence ,  for $n_1+n_2+2 \leq j \leq n_1+n_2+n_3+1$, when we compute the dot product  of the $(n_1+1)$st row of $Q_{{\bf a^1}, {\bf a^2}}$  and the $j$th column of claimed $Q_{{\bf a^1}, {\bf a^2}}^{-1}$ (where we have to take into account $C^{-1}$ at the bottom right block) we just have to add $-(1/\widetilde{p})\; ({\bf v}^C)^T$ (the first row of $C^{-1}$) and hence the dot product will be zero.  

Finally, we can see that, for  $n_1+2 \leq j \leq n_1+n_2+1$, the dot product of the $(n_1+1)$st row of $Q_{{\bf a^1}, {\bf a^2}}$  and the $j$th column of $\widetilde{Q}_{{\bf a^1}, {\bf a^2}}$ 
 is zero. The key point is that these are some multiples of  the dot products of the first row of $\widetilde{B}$ with the last $n_2$ columns of $(\widetilde{B})^{-1}$. Here we again use  the equality~(\ref{eq: abcrows2}). 

As a consequence of our case by case analysis  above,  we can explicitly describe the product $Q_{{\bf a^1}, {\bf a^2}}$ and $\widetilde{Q}_{{\bf a^1}, {\bf a^2}}$ as follows. It is a square matrix of size $n_1+n_2+n_3+1$ such that the first $n_1$ and last $n_3$ rows are zero. The middle rows can be described as the juxtaposition of    three block matrices of sizes $(n_2+1) \times n_1$,  $(n_2+1) \times (n_2+1)$, $(n_2+1) \times n_3$, respectively.  The first row of the first block is $(1/\widetilde{p})\; ({\bf u}^A)^T$ and all other rows are zero. The second block is the $(n_2+1) \times (n_2+1)$ identity matrix and the first row of the third block is 
$(1/\widetilde{p})\; ({\bf v}^C)^T$ and all other rows are zero.

Recall that $\widetilde{Q}_{{\bf a^1}, {\bf a^2}}$ is obtained by deleting $A^{-1}$ and $C^{-1}$ from the matrix depicted in Figure~\ref{fig: inversematrix}. Therefore, based on the previous paragraph and the fact that  $$AA^{-1}=I\;  \mbox{and} \; CC^{-1}=I,$$ we see that the product of the matrix $Q_{{\bf a^1}, {\bf a^2}}$ depicted in Figure~\ref{fig: matrix} and the matrix depicted in Figure~\ref{fig: inversematrix} is equal to the identity matrix. This is because the nonzero row  $(1/\widetilde{p})\; ({\bf u}^A)^T$ in the previous paragraph will be cancelled out by the last row of $A^{-1}$, which is equal to $-(1/\widetilde{p})\; ({\bf u}^A)^T$ and similarly, the nonzero row  $(1/\widetilde{p})\; ({\bf v}^C)^T$ in the previous paragraph will be cancelled out by the first row of $C^{-1}$, which is equal to $-(1/\widetilde{p})\; ({\bf v}^C)^T$. We would like to emphasize that we have also used the fact that $\widetilde{B}(\widetilde{B})^{-1}=I$, while calculating the product of $Q_{{\bf a^1}, {\bf a^2}}$ and $\widetilde{Q}_{{\bf a^1}, {\bf a^2}}$. 
\end{proof} 

\end{document}
